%%%%%%%%%%%%%%%%%%%%%%%%%%%%%%%%%%%%%%%%%%%%%%%%%%%%%%%%%%%%%%%%%%%%%%%%%%%%%%%%
%% PHASE 1 — DETAILED RESEARCH EXECUTION TABLES
%% Topic Selection & Problem Definition (Chapter 1)
%%%%%%%%%%%%%%%%%%%%%%%%%%%%%%%%%%%%%%%%%%%%%%%%%%%%%%%%%%%%%%%%%%%%%%%%%%%%%%%%
\normalsize\normalfont\color{black}

%% ── Table 1: Input / Process / Output ──
Table~\ref{tab:phase1_ipo} maps the eight sequential steps of topic selection and problem definition, specifying the inputs, analytical processes, and deliverables that culminate in seven SMART-validated research objectives.

\needspace{4\baselineskip}
\begingroup\singlespacing\tablefontsize\tablestripes
\begin{xltabular}{\textwidth}{@{} >{\raggedright\arraybackslash}p{0.14\textwidth} X X X @{}}
\caption[Phase~1 --- Input, Process, and Output: Topic Selection]{Phase~1 --- Input, Process, and Output: Topic Selection and Problem Definition}
\label{tab:phase1_ipo}\\
\toprule
\thead{Step} & \thead{Input} & \thead{Process} & \thead{Output} \\
\midrule
\endfirsthead
\endhead
\bottomrule
\endlastfoot
Interest area        & Personal research interest, industry trends                 & Identify broad research theme                  & Research domain candidates \\
\addlinespace
Domain selection     & Academic fields, discipline scope                           & Map interest to academic discipline             & Defined domain: AI governance \\
\addlinespace
Topic narrowing      & Domain literature, industry reports                         & Refine topic scope using [Construct + Outcome + Context] & Specific research topic \\
\addlinespace
Industry problem     & Industry reports, practitioner issues, \$4.5T health burden & Analyse real-world challenges                   & 4 practical problem statements (P1--P4) \\
\addlinespace
Academic gap         & Scholarly articles, literature review                       & Identify missing research                       & 12 research gaps \\
\addlinespace
Problem statement    & Problem + gap                                              & Formulate research problem context               & Formal research problem statement \\
\addlinespace
Research questions   & Problem statement                                          & Convert problem into investigative questions     & 8 research questions (RQ1--RQ8) \\
\addlinespace
Research objectives  & Research questions                                         & Convert questions into measurable objectives     & 8 SMART-validated objectives (O1--O8) \\
\end{xltabular}
\endgroup

\vspace{1.5em}

%% ── Table 2: Quality Validation Framework ──
Table~\ref{tab:phase1_quality} documents the six quality dimensions applied to validate the research design during Phase~1.

\needspace{3\baselineskip}
\begingroup\singlespacing\tablefontsize\tablestripes
\begin{xltabular}{\textwidth}{@{} p{3.2cm} X @{}}
\caption{Phase~1 --- Research Design Quality Framework}
\label{tab:phase1_quality}\\
\toprule
\thead{Dimension} & \thead{Method} \\
\midrule
\endfirsthead
\endhead
\bottomrule
\endlastfoot
Quality check        & Ensure research problem is clear, focused, and researchable \\
\addlinespace
Justification        & Demonstrate importance of problem for academia and industry (\$4.5T health burden, clinical AI trust deficit) \\
\addlinespace
Framework reference  & Research design frameworks, problem statement models, SMART validation \\
\addlinespace
Techniques           & Literature scanning, keyword analysis, gap analysis, topic narrowing formula \\
\addlinespace
Evaluation           & Supervisor feedback, literature coverage review, examiner Q\&A alignment \\
\addlinespace
Benchmarking         & Compare with published DBA/PhD thesis research problems in AI governance domain \\
\end{xltabular}
\endgroup

\vspace{1.5em}

%% ── Table 3: Academic Readiness Evaluation ──
Table~\ref{tab:phase1_readiness} compares PhD, DBA, and professor-level readiness expectations across six evaluation criteria for the problem definition phase.

\needspace{4\baselineskip}
\begingroup\singlespacing\tablefontsize\tablestripes
\begin{xltabular}{\textwidth}{@{} >{\raggedright\arraybackslash}p{0.14\textwidth} >{\raggedright\arraybackslash}p{0.26\textwidth} >{\raggedright\arraybackslash}p{0.28\textwidth} >{\raggedright\arraybackslash}p{0.26\textwidth} @{}}
\caption{Phase~1 --- Professor, PhD, and DBA Readiness Evaluation}
\label{tab:phase1_readiness}\\
\toprule
\thead{Criteria} & \thead{PhD Readiness} & \thead{DBA Readiness} & \thead{Professor Expectation} \\
\midrule
\endfirsthead
\endhead
\bottomrule
\endlastfoot
Problem clarity       & 4 problem statements (P1--P4) ground the \$4.5T health burden and ASD diagnostic fragmentation in IS/AI governance theory & Clinical AI trust deficit quantified: siloed pipelines, no unified governance, ASD domain unaddressed simultaneously & Examiner verifies P1--P4 derive from 12 research gaps (G1--G12) with literature traceability in Section~1.4 \\
\addlinespace
Research contribution & TAM-TOE-RBV integration creates novel mediation model (agentic AI + RAG explainability) absent from prior IS literature & RGAIG framework provides actionable governance roadmap: 525-module taxonomy, 4-level escalation ladder, ROI 135--2{,}717\% & Contributions C1--C8 each map to a gap, hypothesis, and empirical evidence; non-overlapping novelty claims \\
\addlinespace
Literature grounding  & 120+ sources screened via PRISMA; 6 thematic clusters covering DL, agentic AI, RAG, BCI, governance, and trust & Industry reports (WHO, FDA, EU AI Act) alongside 70+ applied studies; \$4.5T burden cited with provenance & Citation density 3--5 per paragraph in Ch.2; $\geq$60\% within 5 years; seminal works (Davis 1989, Hevner 2004) anchored \\
\addlinespace
Research scope        & 5 hypotheses (H1--H5), 8 RQs (RQ1--RQ8), 8 objectives spanning ASD-focused AI diagnostics with mediation and moderation paths & ASD (ASD, PD, epilepsy, stress, cancer) with 19 ASD subjects and 305~GB real clinical data & Scope validated as feasible: DSR + computational experiments on public datasets; no prospective clinical trial needed \\
\addlinespace
Feasibility           & 198 trained models already operational; agenticfinder codebase (Python, joblib, ChromaDB) demonstrates artefact completeness & Real deployment pipeline: Emotiv EPOC X wearable BCI, real-time inference at 97.67\% $\pm$ 2.1\% accuracy, 58~GB raw data processed & Methodology timeline realistic: secondary data from PhysioNet/\allowbreak{}Kaggle/\allowbreak{}OpenNeuro; bootstrap replaces parametric assumptions \\
\addlinespace
Novelty               & To our knowledge, the first unified ASD-focused AI diagnostic framework integrating DL + agentic AI + RAG + governance across ASD & To our knowledge, the first 525-module quality taxonomy spanning 8-phase DL, 13-phase GenAI QA, and 10 governance pillars & Novelty statement explicit in Section~1.8 (C1--C8); gap analysis confirms no prior framework at this scope \\
\end{xltabular}
\endgroup

\vspace{1.5em}

%% ── Table 4: Research Methodology Consideration ──
Table~\ref{tab:phase1_methodology} summarises the three methodological approaches considered and the rationale for adopting a mixed-methods design.

\needspace{3\baselineskip}
\begingroup\singlespacing\tablefontsize\tablestripes
\begin{xltabular}{\textwidth}{@{} p{2.2cm} p{2.2cm} X @{}}
\caption{Phase~1 --- Research Methodology Type Consideration}
\label{tab:phase1_methodology}\\
\toprule
\thead{Method} & \thead{This Study} & \thead{Rationale} \\
\midrule
\endfirsthead
\endhead
\bottomrule
\endlastfoot
Quantitative       & Primary    & Hypothesis testing via DL classification, statistical significance ($p < .05$), bootstrap CI \\
\addlinespace
Qualitative        & Secondary  & Expert panel validation (12-member), thematic coding of RAG outputs, reflexivity statement \\
\addlinespace
Mixed methods      & Adopted    & Quantitative results triangulated with qualitative expert assessment and scenario testing \\
\end{xltabular}
\endgroup

\vspace{1.5em}

%% ── Table 5: Variable Taxonomy (Phase 1 Preview) ──
Table~\ref{tab:phase1_variables} presents a preliminary classification of the research variables, distinguishing the independent, dependent, mediating, moderating, and control constructs.

\needspace{3\baselineskip}
\begingroup\singlespacing\tablefontsize\tablestripes
\begin{xltabular}{\textwidth}{@{} p{3.2cm} p{2.5cm} X @{}}
\caption{Phase~1 --- Research Variable Taxonomy Preview}
\label{tab:phase1_variables}\\
\toprule
\thead{Variable} & \thead{Type} & \thead{Definition} \\
\midrule
\endfirsthead
\endhead
\bottomrule
\endlastfoot
Technology integration (AI/DL)  & Independent variable  & Multi-modal diagnostic models across ASD \\
\addlinespace
Organisational capability          & Mediator (M1)         & Agentic automation, pipeline orchestration, workflow efficiency \\
\addlinespace
Human capability                   & Mediator (M2)         & RAG explainability, clinical trust, expert interpretability \\
\addlinespace
Business value                     & Dependent variable    & Diagnostic accuracy, operational efficiency, governance compliance \\
\addlinespace
Change readiness                   & Moderator             & Digital culture, leadership support, regulatory preparedness \\
\addlinespace
Domain, dataset size, signal type  & Control variables     & Held constant across hypothesis tests \\
\end{xltabular}
\endgroup

\vspace{1.5em}

%% ── Table 6: Evaluation and Benchmarking ──
Table~\ref{tab:phase1_benchmark} enumerates the seven evaluation metrics and benchmarking standards used to assess Phase~1 deliverables.

\needspace{3\baselineskip}
\begingroup\singlespacing\tablefontsize\tablestripes
\begin{xltabular}{\textwidth}{@{} p{3.2cm} X @{}}
\caption{Phase~1 --- Evaluation Metrics and Benchmarking}
\label{tab:phase1_benchmark}\\
\toprule
\thead{Metric / Benchmark} & \thead{Standard} \\
\midrule
\endfirsthead
\endhead
\bottomrule
\endlastfoot
Problem clarity              & Clearly defined in first 2 pages of Chapter~1 \\
\addlinespace
Literature support           & At least 20--30 sources supporting the problem statement \\
\addlinespace
Research gap articulation    & Explicit statement of academic and practical gap \\
\addlinespace
Problem relevance            & Linked to industry challenge (\$4.5T health burden, clinical AI adoption) \\
\addlinespace
Research novelty             & No prior unified ASD-focused AI governance framework for biomedical diagnosis \\
\addlinespace
Objective measurability      & All 8 objectives validated against SMART criteria \\
\addlinespace
Hypothesis testability       & All 5 hypotheses have specified statistical test and acceptance threshold \\
\end{xltabular}
\endgroup

\vspace{1.5em}

%% ── Table 7: Research Evidence Chain ──
Table~\ref{tab:phase1_evidence} reports the eight-link evidence chain connecting the initial research interest to the final set of SMART-validated objectives.

\needspace{3\baselineskip}
\begingroup\singlespacing\tablefontsize\tablestripes
\begin{xltabular}{\textwidth}{@{} p{3.2cm} X @{}}
\caption{Phase~1 --- Research Evidence Chain}
\label{tab:phase1_evidence}\\
\toprule
\thead{Chain Link} & \thead{Evidence} \\
\midrule
\endfirsthead
\endhead
\bottomrule
\endlastfoot
Interest area            & Responsible AI in healthcare enterprises \\
\addlinespace
Research domain          & AI governance / Information Systems / Digital Transformation \\
\addlinespace
Topic                    & Unified AI framework for ASD diagnosis with governance \\
\addlinespace
Practical problem        & ASD-focused fragmentation, clinical trust deficit, \$4.5T health burden \\
\addlinespace
Academic gap             & No unified framework integrating DL, agentic AI, RAG, and governance \\
\addlinespace
Research problem         & 4 problem statements (P1--P4) formally defined \\
\addlinespace
Research questions       & 8 research questions (RQ1--RQ8) aligned to problem statements \\
\addlinespace
Research objectives      & 8 SMART-validated objectives (O1--O8) \\
\end{xltabular}
\endgroup

\vspace{1.5em}

%% ── Table 8: SMART Objective Validation Matrix ──
Table~\ref{tab:phase1_smart} catalogues each of the seven research objectives against the five SMART criteria, demonstrating that every objective is specific, measurable, achievable, relevant, and time-bound.

\needspace{18\baselineskip}
\begingroup\singlespacing\tablefontsize\tablestripes
\begin{xltabular}{\textwidth}{@{} p{0.22\textwidth} X X X X X @{}}
\caption{Phase~1 --- SMART Validation of Research Objectives}
\label{tab:phase1_smart}\\
\toprule
\thead{Objective} & \thead{Specific} & \thead{Measurable} & \thead{Achievable} & \thead{Relevant} & \thead{Time-Bound} \\
\midrule
\endfirsthead
\endhead
\bottomrule
\endlastfoot
O1: Develop unified RGAIG framework & Targets ASD with 5-layer architecture & 525-module taxonomy with pass/fail metrics & DSR artefact using public datasets & Addresses G1: no unified ASD-focused AI governance & Deliverable within DBA timeline (2024--2026) \\
\addlinespace
O2: Validate ASD-focused classification & ASD ensemble accuracy across ASD, PD, epilepsy, stress, cancer & Accuracy, F1, AUC, bootstrap CI & 131 subjects + 20{,}000 images, 198 trained models available & Tests H1 directly & Results reported in Ch.4 \\
\addlinespace
O3: Compare DL vs classical baselines & CNN-LSTM vs RF/\allowbreak{}SVM/\allowbreak{}XGBoost & Paired t-test, McNemar's, Cohen's d & 6 architectures already trained & Tests H2, addresses G3 & Comparison in Section~4.4.3 \\
\addlinespace
O4: Identify discriminative features & Top-k features per domain via SHAP & SHAP values, permutation importance & MNE-Python + scikit-learn pipeline & Tests H3, addresses G4 & Feature analysis in Section~4.4.4 \\
\addlinespace
O5: Evaluate agentic automation & Pipeline orchestration efficiency & Latency, throughput, error rate & agenticfinder operational & Tests H4, addresses G5 & Pipeline metrics in Section~4.4.5 \\
\addlinespace
O6: Assess RAG explainability & Domain-specific RAG explanations & RAGAS scores, expert panel kappa & ChromaDB 1.4~GB embeddings & Tests H5, addresses G6 & RAG evaluation in Section~4.5 \\
\addlinespace
O7: Design governance framework & 10-pillar governance taxonomy & Pass rates per pillar, compliance score & RGAIG 525-module taxonomy & Addresses G7--G12 & Framework in Section~3.6 \\
\end{xltabular}
\endgroup

\vspace{1.5em}

%% ── Table 9: Research Domain Justification ──
Table~\ref{tab:phase1_domain} presents the eight justification criteria that guided the selection of AI governance in biomedical diagnosis as the research domain.

\needspace{4\baselineskip}
\begingroup\singlespacing\tablefontsize\tablestripes
\begin{xltabular}{\textwidth}{@{} p{3.2cm} X @{}}
\caption{Phase~1 --- Research Domain Selection Justification}
\label{tab:phase1_domain}\\
\toprule
\thead{Criterion} & \thead{Justification} \\
\midrule
\endfirsthead
\endhead
\bottomrule
\endlastfoot
Academic relevance      & AI governance is an emerging IS subfield; no unified ASD diagnostic framework exists in prior literature (G1) \\
\addlinespace
Industry significance   & Global healthcare AI market projected \$45.2B by 2030; \$4.5T annual health burden demands AI-driven solutions \\
\addlinespace
Theoretical grounding   & TAM \citep{davis1989tam}, TOE \citep{tornatzky1990processes}, RBV \citep{barney1991firm} provide established theoretical lenses \\
\addlinespace
Methodological fit      & DSR (Hevner, 2004; Peffers, 2007) supports artefact design; mixed methods enable triangulation \\
\addlinespace
Data availability       & Public datasets: ABIDE-II (ASD), PPMI (PD), CHB-MIT (epilepsy), DEAP/SAM40 (stress), OpenNeuro (cancer) \\
\addlinespace
Practical impact        & Framework directly applicable to hospital CIOs/CTOs; ROI 135--2{,}717\% projected; EU AI Act compliance pathway \\
\addlinespace
Novelty                 & First study integrating DL + agentic AI + RAG + governance across ASD simultaneously \\
\addlinespace
DBA suitability         & Combines theoretical contribution (TAM-TOE-RBV integration) with managerial/practical impact (RGAIG deployment roadmap) \\
\end{xltabular}
\endgroup

\vspace{1.5em}

%% ── Table 10: Quality Checklist ──
Table~\ref{tab:phase1_checklist} reports the nine-item quality checklist with pass/fail status and supporting evidence for Phase~1 deliverables.

\needspace{3\baselineskip}
\begingroup\singlespacing\tablefontsize\tablestripes
\begin{xltabular}{\textwidth}{@{} p{0.36\textwidth} >{\centering\arraybackslash}p{0.08\textwidth} p{0.48\textwidth} @{}}
\caption{Phase~1 --- Quality Checklist}
\label{tab:phase1_checklist}\\
\toprule
\thead{Checklist Item} & \thead{Status} & \thead{Evidence} \\
\midrule
\endfirsthead
\endhead
\bottomrule
\endlastfoot
Is the research problem clearly defined?                 & $\checkmark$ & 4 problem statements (P1--P4) defined in Section~1.2; formal problem statement in Section~1.3 \\
\addlinespace
Is the problem supported by literature?                  & $\checkmark$ & 30+ sources cited in Sections~1.1--1.3; \$4.5T global health burden referenced with WHO data \\
\addlinespace
Is the research gap identified?                          & $\checkmark$ & 12 research gaps (G1--G12) enumerated in Section~1.4 with literature traceability \\
\addlinespace
Is the problem relevant to the domain?                   & $\checkmark$ & Clinical AI trust deficit and ASD-focused fragmentation linked to industry reports in Section~1.1 \\
\addlinespace
Is the research feasible?                                & $\checkmark$ & 5 public datasets identified; 19 ASD subjects; real codebase operational (Section~1.6) \\
\addlinespace
Are research questions aligned with the problem?         & $\checkmark$ & 8 research questions (RQ1--RQ8) mapped to P1--P4 in Table~\ref{tab:phase1_evidence} \\
\addlinespace
Are objectives SMART-validated?                          & $\checkmark$ & 8 objectives (O1--O8) with SMART validation matrix in Section~1.5 \\
\addlinespace
Is the theoretical framework identified?                 & $\checkmark$ & TAM-TOE-RBV integration introduced in Section~1.7; detailed in Chapter~2 \\
\addlinespace
Is the research novelty explicitly stated?               & $\checkmark$ & Contributions C1--C8 listed in Section~1.8; no prior unified ASD-focused AI governance framework identified \\
\end{xltabular}
\endgroup

\clearpage
